% Part: model-theory
% Chapter: models-of-arithmetic
% Section: computable-models

\documentclass[../../../include/open-logic-section]{subfiles}

\begin{document}

\olfileid{mod}{mar}{cmp}
\section{পাটিগণিতের গণনসাধ্য মডেল}

\begin{explain}
মানক মডেল~$\Struct{N}$-এর দুটি সুন্দর বৈশিষ্ট্য আছে। তার সংজ্ঞাক্ষেত্র
স্বাভাবিক সংখ্যা~$\Nat$; অর্থাৎ পাটিগণিতের ভাষায় আমরা যে ধরনের বস্তু
নিয়ে কথা বলতে চাই, তার উপাদানগুলি ঠিক সেই বস্তু, এবং মানক
সংখ্যাপদ~$\num{n}$ সত্যিই~$n$-কেই নির্দেশ করে। অন্য সুন্দর বৈশিষ্ট্যটি
হলো, $\Lang{L_A}$-এর অ-যৌক্তিক সংকেতগুলির সব ব্যাখ্যাই
\emph{গণনসাধ্য}। $\Assign{\prime}{N}$, $\Assign{+}{N}$ ও
$\Assign{\times}{N}$ হিসেবে কাজ করা উত্তরসূরি, যোগ ও গুণ অপেক্ষকগুলি
সংখ্যার গণনসাধ্য অপেক্ষক। (যেমন, টুরিং যন্ত্র দিয়ে গণনসাধ্য, অথবা
আদিম পুনরাবৃত্তি দিয়ে সংজ্ঞেয়।) আর $\Struct{N}$-এর ক্ষুদ্রতর-সম্বন্ধ,
অর্থাৎ $\Assign{<}{N}$, নির্ণেয়।

$\Th{Q}$ ও~$\Th{PA}$-র মতো পাটিগাণিতিক তত্ত্বের অমানক মডেলে অমানক
উপাদান থাকতেই হয়। তাই তাদের সংজ্ঞাক্ষেত্রে সাধারণত~$\Nat$-এর সঙ্গে
অতিরিক্ত !!{element}s থাকে। কিন্তু যেকোনো গণনীয় !!{structure} যেকোনো
!!{denumerable} সেটের উপর নির্মাণ করা যায়, বিশেষত~$\Nat$-এর উপরও।
তাই $\Nat$ সংজ্ঞাক্ষেত্রবিশিষ্ট অমানক মডেলও আছে। অবশ্য এমন মডেল
$\Struct{M}$-এ অন্তত কয়েকটি সংখ্যা তাদের স্বাভাবিক ভূমিকা পালন করতে
পারে না; কারণ এমন কোনো~$k$ থাকতেই হবে, যা সব~$n \in \Nat$-এর জন্য
$\Value{\num{n}}{M}$ থেকে ভিন্ন।
\end{explain}

\begin{defn}
$\Lang{L_A}$-এর কোনো !!^a{structure}~$\Struct{M}$ ঠিক তখনই
\emph{গণনসাধ্য}, যখন $\Domain{M}=\Nat$; $\Assign{\prime}{M}$,
$\Assign{+}{M}$ ও~$\Assign{\times}{M}$ গণনসাধ্য অপেক্ষক; এবং
$\Assign{<}{M}$ একটি নির্ণেয় সম্বন্ধ।
\end{defn}

\begin{ex}\ollabel{ex:comp-model-q}
\olref[mdq]{ex:model-K-of-Q}-এর গঠন~$\Struct{K}$ মনে করো। তার
সংজ্ঞাক্ষেত্র ছিল $\Domain{K}=\Nat \cup \{a\}$ এবং ব্যাখ্যাগুলি ছিল
\begin{align*}
  \Assign{\Obj{0}}{K} & = 0\\
  \Assign{\prime}{K}(x) & =
  \begin{cases}
    x+1 & \text{যদি $x\in \Nat$}\\
    a & \text{যদি $x = a$}
  \end{cases}\\
  \Assign{+}{K}(x, y) & =
  \begin{cases}
    x+y & \text{যদি $x$, $y \in\Nat$}\\
    a & \text{অন্যথায়}
  \end{cases}\\
  \Assign{\times}{K}(x, y) & =
  \begin{cases}
    xy & \text{যদি $x$, $y \in\Nat$}\\
    0 & \text{যদি $x=0$ অথবা $y=0$}\\
    a & \text{অন্যথায়}\\
  \end{cases}\\
  \Assign{<}{K} & =
  \Setabs{\tuple{x,y}}{x, y \in \Nat \text{ এবং } x<y} \cup
  \Setabs{\tuple{x,a}}{x \in \Domain{K}}
\end{align*}
% BN-SRC-132 TARGET-CORRECTION-BEGIN
\par\smallskip\noindent\textnormal{\small\textbf{উৎস-সংশোধন (BN-SRC-132):} হিমায়িত ইংরেজি উৎসে শেষ সেট-নির্মাণের সাজানো জোড়ে প্রথম স্থানটি~$x$, কিন্তু শর্তে অসংজ্ঞায়িত~$n\in\Domain{K}$ লেখা হয়েছে। বাংলা সূত্রে বাঁধা চলরাশির সঙ্গে মিলিয়ে শর্তটি $x\in\Domain{K}$ করা হয়েছে।} % BN-SRC-132 TARGET-CORRECTION-END
কিন্তু $\Domain{K}$ !!{denumerable}, তাই সেটি~$\Nat$-এর সঙ্গে সমসংখ্যক।
যেমন, $g\colon \Nat \to \Domain{K}$, যেখানে $g(0)=a$ এবং $n>0$ হলে
$g(n)=n+1$, একটি !!a{bijection}। একে $\Th{Q}$-এর নতুন
মডেল~$\Struct{K'}$ ও~$\Struct{K}$-এর মধ্যকার সমরূপতায় বদলাতে পারি।
$\Struct{K'}$-এ $\Lang{L_A}$-এর সংকেতগুলিকে ভিন্ন
অপেক্ষক ও সম্বন্ধ দিতে হবে, কারণ~$\Nat$-এর ভিন্ন !!{element}s-গুলি
মানক ও অমানক সংখ্যার ভূমিকা পালন করে।

নির্দিষ্টভাবে, $0$ এখন ক্ষুদ্রতম মানক সংখ্যার নয়, $a$-এর ভূমিকা পালন
করে। ক্ষুদ্রতম মানক সংখ্যা এখন~$1$। তাই
$\Assign{\Obj{0}}{K'}=1$ নিযুক্ত করি। উত্তরসূরি অপেক্ষকটিও এখন আলাদা:
কোনো মানক সংখ্যা, অর্থাৎ $n>0$, দিলে এখনও $n+1$ ফেরায়। কিন্তু $0$ এখন
$a$-এর ভূমিকা পালন করে, আর সেটি নিজেরই উত্তরসূরি। তাই
$\Assign{\prime}{K'}(0)=0$। একইভাবে যোগ ও গুণের জন্য পাই
\begin{align*}
\Assign{+}{K'}(x, y) & =
  \begin{cases}
    x+y-1 & \text{যদি $x$, $y >0$}\\
    0 & \text{অন্যথায়}
  \end{cases}\\
  \Assign{\times}{K'}(x, y) & =
  \begin{cases}
    1 & \text{যদি $x = 1$ অথবা $y = 1$}\\
    xy - x - y + 2 & \text{যদি $x$, $y > 1$}\\
    0 & \text{অন্যথায়}\\
  \end{cases}
\end{align*}
আর $\tuple{x,y}\in\Assign{<}{K'}$ যদি এবং কেবল যদি $x<y$, $x>0$ ও
$y>0$; অথবা $y=0$।

এই অপেক্ষকগুলি সবই স্বাভাবিক সংখ্যার গণনসাধ্য অপেক্ষক এবং
$\Assign{<}{K'}$, $\Nat$-এর উপর একটি নির্ণেয় সম্বন্ধ---কিন্তু তারা
$\Nat$-এর উত্তরসূরি, যোগ ও গুণ অপেক্ষকগুলির সমান নয়, এবং
$\Assign{<}{K'}$-ও $\Nat$-এর উপর~$<$ সম্বন্ধটির সমান নয়।
\end{ex}

\begin{prob}
$\Domain{L'}=\Nat$ এমন একটি !!a{structure}~$\Struct{L'}$ দাও, যা
\olref[mod][mar][mdq]{ex:model-L-of-Q}-এর $\Struct{L}$-এর সমরূপ।
\end{prob}

\begin{explain}
\Olref{ex:comp-model-q} দেখায়, $\Th{Q}$-এর $\Nat$ সংজ্ঞাক্ষেত্রবিশিষ্ট
গণনসাধ্য অমানক মডেল আছে। কিন্তু নিচের ফলটি দেখায়, $\Th{PA}$-র মডেলের
ক্ষেত্রে কথাটি সত্য নয় (এবং তাই $\Th{TA}$-র মডেলের ক্ষেত্রেও নয়)।
\end{explain}

\begin{thm}[টেনেনবাউমের উপপাদ্য]
$\Struct{N}$ সমরূপতা পর্যন্ত $\Th{PA}$-র একমাত্র গণনসাধ্য মডেল।
\end{thm}
% BN-SRC-133 TARGET-CORRECTION-BEGIN
\par\smallskip\noindent\textnormal{\small\textbf{উৎস-সংশোধন (BN-SRC-133):} হিমায়িত বিবৃতিতে $\Struct{N}$-কে আক্ষরিকভাবে একমাত্র গণনসাধ্য মডেল বলা হয়েছে। কিন্তু $\Nat$-এর কোনো গণনসাধ্য বিজেকশনের মাধ্যমে ব্যাখ্যাগুলি স্থানান্তর করলে একই মানক মডেলের ভিন্ন গণনসাধ্য উপস্থাপন পাওয়া যায়। টেনেনবাউমের উপপাদ্যের সঠিক সিদ্ধান্ত হলো, কোনো গণনসাধ্য অমানক মডেল নেই; তাই বাংলা পাঠে ``সমরূপতা পর্যন্ত'' যোগ করা হয়েছে।} % BN-SRC-133 TARGET-CORRECTION-END
  
\end{document}
